\section{Random tail covariance matrices}\label{app:gram}
Let $Q$ be the covariance kernel retaining only degrees greater than an
odd cutoff $r$, and write $\tau=\tau_r$, $\ell=\lambda_{r+1}$, and
$s_2=s_{2,r}$. This section proves the sample event needed later, without
assuming that a Gaussian approximation remains valid after conditioning.
The entire argument concerns the exact limiting Gaussian covariance.

\begin{lemma}[Far-pair decay]\label{lem:far}
There are constants $c_0,C>0$ such that, when $r\le D/8$ is odd,
\[
 \max_{d/4\le\dist(x,y)\le3d/4}|Q(x,y)|
 \le\eta_r:=C\sqrt D\,e^{-r/8}+Ce^{-c_0D}.
\]
\end{lemma}
\begin{proof}
The full-cube Krawtchouk polynomial satisfies
\[
 \sum_{j=0}^D\mathsf K_j(h;D)z^j=(1+z)^{D-h}(1-z)^h.
\]
If $\delta D\le h\le(1-\delta)D$, put $b=\min(h,D-h)$.
On $|z|=t$ the generating function is bounded in absolute value by
\[
 (1+t^2)^b(1+t)^{D-2b}
 \le(1+t)^D\exp\!\left[-\frac{2\delta Dt}{(1+t)^2}\right].
\]
Cauchy's coefficient bound at $t=j/(D-j)$ and the uniform Stirling
lower bound for $\binom Dj$ therefore give, for $1\le j<D$,
\begin{equation}\label{eq:kraw}
 \frac{|\mathsf K_j(h;D)|}{\binom Dj}
 \le C\sqrt D\exp[-2\delta j(1-j/D)].
\end{equation}
For reduced-cube points at distance in $[d/4,3d/4]$, their augmented
full-cube distance lies in $[D/8,7D/8]$ for large $d$.
Because $r$ is odd, the exact tail is
$Q(x,y)=\sum_{j>r,\ j\text{ odd}}a_j\mathsf K_j(h;D)$.
For $r<j\le D-r$, \eqref{eq:kraw} gives a bound
$C\sqrt D e^{-r/8}$ after summation, since the total spectral mass
is one. For $j>D-r$, monotonicity and an odd cutoff $b_D$ near $D/4$
give
\[
 \sum_{j>D-r}\binom Dj a_j
 \le a_{b_D}\sum_{k<r}\binom Dk
 \le C D^{-\beta-1/2}
       \exp[-D(h(1/4)-h(1/8))+O(1)].
\]
Here $h(t)=-t\log t-(1-t)\log(1-t)$ denotes the binary entropy in nats.
The positive entropy gap absorbs the polynomial factor and proves the
claim. Adjust constants for bounded $d$.
\end{proof}

\begin{lemma}[A tail Gram deviation bound]\label{lem:gram-general}
Let $X_1,\ldots,X_n$ be independent uniform cube points, with
$1\ll n\le M/4$. Let
\[
 B_{\mathrm{near}}=\#\{y:\dist(x,y)\notin[d/4,3d/4]\}
 \le 2e^{d h(1/4)},\qquad q=2\log M.
\]
The count is independent of $x$. For every $z>0$,
\begin{align}
 &\mathbb P\{\|Q_{SS}-\tau I_n\|_{\mathrm{op}}>z\tau\}
 \notag\\
 &\quad\le\frac C{z^2}
 \left[\frac{q\eta_r}{\tau}
       +\frac{\sqrt{nq s_2}}{\tau}
       +\frac{n\ell}{\tau}
       +\frac{n(1+B_{\mathrm{near}})}M\right]^2
       +\frac{n^2(1+B_{\mathrm{near}})}M.
 \label{eq:gram-bound}
\end{align}
Here the sampled matrix includes repeated training locations.
\end{lemma}
\begin{proof}
Let $J=Q-\tau I_M$, a zero-diagonal Hermitian matrix. Form a deterministic
mask $J_{\mathrm{far}}$ retaining only pairs at distances in
$[d/4,3d/4]$. Since $Q$ is positive semidefinite with constant diagonal
$\tau$, every entry has absolute value at most $\tau$. Thus
\begin{align*}
 \|J_{\mathrm{far}}\|_{\max}&\le\eta_r,\\
 \|J_{\mathrm{far}}\|_{1\to2}^2&\le M s_2,\\
 \|J_{\mathrm{far}}\|_{\mathrm{op}}
 &\le M\ell+\tau+\tau B_{\mathrm{near}}.
\end{align*}
For the second inequality, Parseval gives squared column norm
$M s_2-\tau^2$ for $J$, and entrywise masking only decreases it.
For the third, $\|J\|\le M\ell+\tau$, while the removed matrix has
maximum absolute row sum at most $\tau B_{\mathrm{near}}$.
The entropy bound on $B_{\mathrm{near}}$ is the elementary binomial
tail bound on distances to a fixed cube point.

Let $R$ retain each of the $M$ coordinates independently with
probability $p=2n/M$. The zero-diagonal case of Tropp's random
principal-submatrix theorem \cite[Theorem 1.1]{tropp2008}, at moment
order two and $q=2\log M$, states
\[
 (\mathbb E\|RJR\|^2)^{1/2}
 \le C\left[q(\mathbb E\|RJR\|_{\max}^2)^{1/2}
       +\sqrt{pq}(\mathbb E\|JR\|_{1\to2}^2)^{1/2}
       +p\|J\|\right].
\]
Apply it to $J_{\mathrm{far}}$ and substitute the preceding bounds:
\begin{equation}\label{eq:bernoulli-gram}
 (\mathbb E\|R J_{\mathrm{far}}R\|^2)^{1/2}
 \le C[q\eta_r+\sqrt{nq s_2}+n\ell
               +(n/M)\tau(1+B_{\mathrm{near}})].
\end{equation}
The logarithm here is that of the original dimension $M$, not the
sample size.

To pass to a uniform subset of size $n$, use a uniform permutation of
all cube points, independently of $B\sim\operatorname{Bin}(M,2n/M)$.
Its first $B$ elements form the Bernoulli subset. On $B\ge n$, the
first $n$ principal submatrix is a compression of the larger one, so
its operator norm cannot be larger. The size-$n$ event is independent
of $B$, and $\mathbb P(B\ge n)\ge1/2$ for large $n$. Thus Markov's
inequality applied to \eqref{eq:bernoulli-gram} gives the same fixed-size
probability bound with a constant factor loss.

Conditional on no collisions, an independent size-$n$ sample is a
uniform size-$n$ subset, up to an irrelevant ordering. Collisions have
probability at most $n(n-1)/(2M)$. The probability of any sampled pair
violating the distance condition is at most $n^2B_{\mathrm{near}}/M$.
On their complement, the sampled matrices $J$ and $J_{\mathrm{far}}$
agree. Adding these exceptional probabilities proves
\eqref{eq:gram-bound}.
\end{proof}

